\section{What Is Established and What Remains Open}

Generative modeling combines mature mathematics with rapidly evolving empirical practice.
It is therefore especially important to distinguish settled ideas from open questions.

\subsection*{What is well established}

Several core principles are by now firmly understood.

\paragraph{EM as lower-bound optimization.}
The EM framework provides a principled way to optimize latent-variable likelihoods through alternating inference and parameter updates.

\paragraph{ELBO logic in variational learning.}
The evidence lower bound gives a clean variational foundation for approximate latent-variable inference in deep models.

\paragraph{Change of variables in flows.}
The flow framework is mathematically exact once one specifies an invertible transformation with tractable Jacobian determinant.

\paragraph{Diffusion forward and reverse formalism.}
The Gaussian forward process, its closed-form marginals, and reverse-process approximation framework are well established.

\subsection*{What remains difficult}

Other aspects are still less resolved.

\paragraph{GAN stability and mode coverage.}
Adversarial training remains delicate, and ensuring both realism and diversity is challenging.

\paragraph{Semantic controllability in latent spaces.}
Latent-variable models often learn useful structure, but guaranteeing interpretable and controllable semantics remains difficult.

\paragraph{Why diffusion scales so well perceptually.}
Diffusion models are empirically powerful, but the full explanation for their perceptual success across domains is still incomplete.

\paragraph{Representation versus generation tradeoffs.}
How sample quality, latent usefulness, controllability, and likelihood should be balanced remains an open modeling question.

\paragraph{Evaluation across objectives.}
Comparing generative models is inherently difficult because they optimize different goals:
likelihood, realism, diversity, reconstruction, or denoising quality may not align.

\subsection*{Closing perspective}

This chapter has organized generative modeling around four major routes:
\begin{itemize}
    \item latent-variable inference,
    \item adversarial learning,
    \item invertible transport,
    \item iterative denoising.
\end{itemize}

Each route offers a different answer to the problem of modeling complex data distributions.
No single paradigm dominates every criterion.

\medskip

The broader lesson is that generative modeling is not a peripheral topic in modern AI.
It is one of the main ways of learning structure, representation, and uncertainty in rich data.

\medskip

The next chapter steps back from specific model classes and asks a more reflective question:

\begin{quote}
what do we actually understand about deep learning, where does it remain unreliable, and what remains genuinely open?
\end{quote}